\documentclass[11pt]{article}
\usepackage{fontspec}
\usepackage[french,english]{babel}
\usepackage{amsmath}
\usepackage{amssymb}
\usepackage{amsthm}
\usepackage{mathrsfs}
\usepackage{tikz-cd}
\usepackage{enumitem}
\usetikzlibrary{arrows.meta}

\newcommand{\printedpage}[1]{\marginpar{\footnotesize p.~#1}}
\newcommand{\msp}[1]{\marginpar{\footnotesize [#1]}}
\newcommand{\cF}{\mathcal{F}}
\newcommand{\cM}{\mathcal{M}}
\newcommand{\Cat}{\mathrm{Cat}}
\newcommand{\Fib}{\mathrm{Fib}}
\newcommand{\Fibf}{\mathrm{Fibf}}
\newcommand{\Parf}{\mathrm{Parf}}
\newcommand{\Fl}{\mathrm{Fl}}
\newcommand{\Comp}{\mathrm{Comp}}
\newcommand{\Catfib}{\mathrm{Catfib}}
\newcommand{\dottedimplies}{\mathrel{\tikz[baseline=-0.6ex]{\draw[dotted, double, double distance=1.5pt, -{Implies}] (0,0) -- (0.9,0);}}}

\begin{document}

\printedpage{24}
% head: A. Grothendieck : Dérivateurs, Ch. XVI, Édition des Universités de Montpellier II et Paris VII
% foot: Version 18 mars 2019

\begin{itemize}[leftmargin=4em]
\item[F)] \footnote[26]{added after the fact.} Cat has sufficiently many $W$-fibrations, \textit{i.e.} Cat equipped with $(W, \Fib_W)$ is a category with cofibrations (proper, and even a Brown one, \textit{i.e.} $X \longrightarrow e$ is in $\Fib_W$ for every $X$ in Cat. % [sic] (propre, et même de Brown, -> parenthesis never closed, closing ")" apparently missing
\end{itemize}

It is clear that c is the special case of b obtained by taking $Y = S$, $q = \mathrm{id}_S$, so b $\Longrightarrow$ c. The implication a $\Longrightarrow$ b is less evident, it is established first in XII, and we will return to it below (4.4.2 below). The rest of the present paragraph is devoted to studying various variants of the preceding conditions a) b) c) and the implications between them.

\msp{page 40}

4.4. Let, more generally,
\[
\cF \subset \Fl(\Cat)
\]
be a set of arrows in Cat, thinking more particularly of the case where $\cF$ is one of the seven sets in diagram (4.2.1) (p. 30) other than $\Fib_W$. We can formulate three properties for $\cF$ similar to those formulated in Loc (7 a b c) in the case of $\cF_W$, let us denote these properties
\[
\mathrm{A}(\cF)\,,\ \mathrm{B}(\cF)\,,\ \mathrm{C}(\cF)\,,\ \text{\footnotemark[27]}
\]
\footnotetext[27]{\textbf{N.B.} These conditions make sense if $\cF \subset \Fl(\cM)$, if $(\cM, W)$ is an arbitrary category of models.}
thus
\[
\left\{
\begin{array}{lcl}
\mathrm{A}(\cF) & : & \cF \subset \Fib_W. \\[1ex]
\mathrm{B}(\cF) & : & \Fl(\cF_S) \cap W_S \subset W_S^{\mathrm{u}} \\
 & & \text{(where $\cF_S$ denotes the full subcategory of $\Cat/S$} \\
 & & \text{consisting of the $f : X \longrightarrow S$ with $f \in \cF$).} \\[1ex]
\mathrm{C}(\cF) & : & \cF \cap W \subset W^{\mathrm{u}}.
\end{array}
\right.
\tag*{(4.4.1)}
\]
It is clear that we have $\mathrm{B}(\cF) \Longrightarrow \mathrm{C}(\cF)$, and I am going to give a condition on $\cF$ which implies that we also have $\mathrm{A}(\cF) \Longrightarrow \mathrm{B}(\cF)$.

\msp{page 41}

(We will apply it in the following cases: % [sic] (Nous -> opening parenthesis never closed in this paragraph, closing ")" apparently missing
\[
\cF \;=\; \cF_W,\ \Fibf_W,\ \Parf_W,\ \cF_{W_0},\ \Fibf_{W_0}.
\tag*{(4.4.2)}
\]
I point out moreover that we will see below (cor. 4.4.10) that
\[
\mathrm{A}(\cF_W) \iff \mathrm{A}(\Fibf_W), \qquad \mathrm{A}(\cF_{W_0}) \iff \mathrm{A}(\Fibf_{W_0}),
\tag*{(4.4.3)}
\]
and likewise for the variants B, C, so that for the implication we want to establish, there is no need to distinguish the cases of $\cF_W$ and $\Fibf_W$, of $\cF_{W_0}$ and $\Fibf_{W_0}$. Moreover it will suffice to look at the sets $\cF_{W_0}$, $\Fibf_{W_0}$, $\Parf_{W_0}$, since the case of the sets relative to $W$ is a special case of these (taking $W_0 = W$). \footnote[28]{This will thus make a priori $3 \times 6 = 18$ different conditions, but because of (4.4.3) there are (up to equivalence) only $4 \times 3 = 12$, and in fact two among these four triples of conditions reduce to a single one (for $\cF = \cF_W$ and $\cF = \cF_{W_0}$), so that in all only \textit{eight} different conditions remain, which we summarize in the summary diagram (4.4.7.1) on page 48. There are even only \textit{four,} $\mathrm{A}_{W_0}$, $\mathrm{B}_{W_0}$, $\mathrm{C}_{W_0}$, $\mathrm{D}_{W_0}$, as functions of $W_0$, the other four A, B, C, D are obtained by specializing to the case where $W_0 = W$.}

\printedpage{25}
% head: A. Grothendieck : Dérivateurs, Ch. XVI, Édition des Universités de Montpellier II et Paris VII
% foot: Version 18 mars 2019

\textbf{Lemma 4.4.1.} \footnote[29]{Holds in an arbitrary category of models.} \textit{Suppose that $\cF$ satisfies the following condition: every arrow $f$ of} Cat \textit{factors as $f = pi$, with $i \in W$, $f \in \cF$. Then the condition} $\mathrm{A}(\cF)$ \textit{on $\cF$ implies that in} Cat \textit{there are ``enough $W$-fibrations'', and consequently} (\textit{by a formal argument well} % [sic] f \in \cF -> p \in \cF, the factorisation is f = pi and lemme 4.4.3 3) restates this condition with p \in \cF
\msp{page 42}
\textit{known}), \textit{that we have} $\mathrm{B}(\Fib_W)$, \textit{and a fortiori} (\textit{since by hypothesis} $\cF \subset \Fib_W$) $\mathrm{B}(\cF)$.

\textbf{N.B.} Indeed it is clear that if we have $\cF \subset \cF'$, then
\[
\mathrm{A}(\cF') \Longrightarrow \mathrm{A}(\cF) \qquad \text{and likewise for B, C.}
\tag*{$(*)$}
\]
Now the preliminary condition on $\cF$ is indeed satisfied in the case of a set of one of the forms $\cF \;=\; \cF_W$, $\Fibf_W$, $\Parf_W$. It suffices to see this in the case of the subsets $\cF_W$, $\Parf_W$ of $\Fibf_W$, \textit{i.e.} that every arrow $f$ of Cat factors as
\[
f = pi,
\]
with $i \in W$, and $p$ being, as desired, in $\Parf_W$, or in $\cF_W$. Now in the case of $\Parf_W$, \marginpar{$\|$}this is what is established in VII (Categories of paths (2)) thanks to the formalism of the $\underline{\mathrm{Ch}}(S) \longrightarrow S \times S$. To find a factorization with $f = qj$, $j \in W$, $q \in \cF_W$, it suffices to see that $p$ factors in this way as $p = q_0 j_0$, we will have % [sic] établie -> établi, agreement with "ce qui"
\[
f = \underbrace{q_0}_{\in \cF_{W_0}} \underbrace{(j_0 i)}_{\in W},
\]
so we can assume $f \in \Parf_W$. But if $X$ is $W$-perfect over $S$, hence a
\msp{page 43}
weak $W$-fibered object over $S$, then in the factorization
\begin{center}
\begin{tikzcd}[column sep=small]
X \arrow[rr, "i"] \arrow[dr, "f"'] & & \Phi_S(X) \arrow[dl, "p"] \\
 & S &
\end{tikzcd}
\end{center}
with $i \in W$ and $p \in \Catfib$ (\textit{cf.} VI), we know that $p$ is also a weak $W$-fibration, \textit{i.e.} $p \in \cF_W$, whence the conclusion. Thus we find the

\textbf{Corollary 4.4.2.} \textit{Suppose that $\cF$ is one of the six sets} (4.4.2) (i.e. \textit{one of the sets in diagram} (4.2.1), \textit{excluding $\mathit{Fib}_W$ and $\mathit{Fib}_{W_0}$}). \textit{Then we have}
\[
\mathrm{A}(\cF) \;\Longrightarrow\; \mathrm{B}(\cF) \;\Longrightarrow\; \mathrm{C}(\cF)\,.
\]

\textbf{Lemma 4.4.3.} \textit{Let $(\cM, W)$ be a category of models, and let $\cF \subset \Fl(\cM)$ satisfy the following conditions}
\begin{itemize}[leftmargin=3em]
\item[1)] \textit{$\cF$ stable under base change, composition, contains the isomorphisms.}
\item[2)] \textit{$\cM$ has a final object, and for every $X \in \mathrm{Ob}\,\cM$, $X \longrightarrow e$ is in $\cF$.}
\item[3)] \textit{Every arrow $f$ of $\cM$ factors as $f = pi$, with $i \in W$, $p \in \cF$} (\textit{this is the condition of lemma} 4.4.1).
\end{itemize}

\printedpage{26}
% head: A. Grothendieck : Dérivateurs, Ch. XVI, Édition des Universités de Montpellier II et Paris VII
% foot: Version 18 mars 2019
\noindent\textit{Then the conditions} $\mathrm{A}(\cF)$, $\mathrm{B}(\cF)$, $\mathrm{C}(\cF)$ \textit{on $\cF$ are equivalent.}

Indeed, by 4.4.1, it remains to prove the implication
\[
\mathrm{C}(\cF) \;\Longrightarrow\; \mathrm{A}(\cF),
\]

\msp{page 44}

\textit{i.e.} this: if $\cF$ also satisfies the condition
\begin{itemize}[leftmargin=3em]
\item[4)] $\cF \cap W \subset W^{\mathrm{u}}$ (or what amounts to the same thing, since $\cF$ is stable under base change, that $\cF \cap W$ is stable under base change),
\end{itemize}
then we conclude from this
\[
\cF \subset \Fib_W.
\]
Now this is none other than the lemma of K. S. \textsc{Brown}, in the form dual to the one given above (3.2.1, page 20).

Among the six sets of arrows considered in (4.4.2), the hypotheses 1) 2) 3) are satisfied for $\Parf_W$ and $\Parf_{W_0}$, to the exclusion of the other four. (The others satisfy all the conditions, with the sole exception of stability under composition \footnote[30]{\textbf{N.B.} For $\Fib_{W_0}$, it is the conditions 1) 2) that hold, and on the other hand, condition 3) that is a problem. Let us denote it by $\mathrm{F}(W_0)$. If it is satisfied (\textit{i.e.} if Cat has enough $W_0$-fibrations), then we also have equivalence of the conditions $\mathrm{A}(\Fib_{W_0})$, $\mathrm{B}(\Fib_{W_0})$, $\mathrm{C}(\Fib_{W_0})$.}). We thus find:

\textbf{Corollary 4.4.4.} \textit{The conditions} $\mathrm{A}(\Parf_W)$, $\mathrm{B}(\Parf_W)$, $\mathrm{C}(\Parf_W)$ \textit{are equivalent} (we denote this condition on $W$ by D). \textit{Likewise the conditions} $\mathrm{A}(\Parf_{W_0})$, $\mathrm{B}(\Parf_{W_0})$, $\mathrm{C}(\Parf_{W_0})$ \textit{are equivalent.}

\msp{page 45}

(We will denote the condition obtained on the pair $W \supset W_0$ of basic localizers by $D_{W_0}$.)

Thus we have (since $\Parf_{W_0} \subset \Parf_W$)
\[
\mathrm{D} = \mathrm{D}_W \;\Longrightarrow\; \mathrm{D}_{W_0}\,,
\tag*{(4.4.4)}
\]
and if $W'_0 \subset W_0$, we will have
\[
\mathrm{D}_{W_0} \;\Longrightarrow\; \mathrm{D}_{W'_0}\,.
\tag*{(4.4.5)}
\]

\textbf{Corollary 4.4.5.} \textit{Let us denote by} A, B, C \textit{respectively the conditions} Loc (7 a), Loc (7 b), Loc (7 c). \textit{Then we have the implications}
\[
\mathrm{A} \;\Longrightarrow\; \mathrm{B} \;\Longrightarrow\; \mathrm{C}\,.
\tag*{(4.4.5.1)}
\]
\textit{Moreover, if we denote} (\textit{for a set $\cF$ of arrows of} Cat) \textit{by} $\Comp(\cF)$ \textit{the condition on $\cF$ that $\cF$ be stable under composition, we have the equivalences}
\[
\mathrm{A} \iff (\mathrm{C} + \Comp(\cF_W)) \iff (\mathrm{C} + \Comp(\Fibf_W)).
\tag*{(4.4.5.2)}
\]

\printedpage{27}
% head: A. Grothendieck : Dérivateurs, Ch. XVI, Édition des Universités de Montpellier II et Paris VII
% foot: Version 18 mars 2019

\textsc{Proof}. The first assertion is contained in corollary 4.4.2, applied to the case $\cF = \cF_W$. On the other hand, if $\cF_W$ satisfies $\Comp(\cF_W)$, it satisfies conditions 1) 2) 3) of Brown's lemma 4.4.3, so combined with C, it implies A. It remains, for the first equivalence, to prove that
\[
A \;\Longrightarrow\; \Comp(\cF_W).
\]

\msp{page 46}

But by definition we have
\[
\cF_W = \Catfib \cap \Fibf_W
\tag*{$(*)$}
\]
and we know that Catfib is stable under composition. So $\cF_W$ is too, provided that the composite $gf$ of two arrows in $\cF_W$ is in $\Fibf_W$. Now by hypothesis A, which means $\cF_W \subset \Fib_W$, we have $f, g \in \Fib_W$, hence $gf \in \Fib_W$, since $\Fib_W$ is stable under composition, a fortiori $gf \in \Fibf_W$.

On the other hand, the preceding formula $(*)$ shows that $\Comp(\Fibf_W) \Longrightarrow \Comp(\cF_W)$, so we get for free the implication $\Longleftarrow$ in the second equivalence (4.4.5.2). The implication in the opposite direction then amounts to
\[
\mathrm{A} \;\Longrightarrow\; \Comp(\Fibf_W),
\]

\msp{page 47}

which will be evident thanks to the result (given below)
\[
\mathrm{A} = \mathrm{A}(\cF_W) \iff \mathrm{A}(\Fibf_W) \qquad \text{and likewise for B, C,}
\]
which thus means that
\[
\cF_W \subset \Fib_W \qquad \text{implies} \qquad \Fibf_W \subset \Fib_W \qquad (\text{so } \Fibf_W = \Fib_W),
\]
now $\Fib_W$ being stable under composition, the same is true of $\Fibf_W$, q.e.d.

\textbf{Corollary 4.4.6.} \textit{Suppose that $\cF_{W_0}$ is stable under composition, which is the case in particular} (\textit{by} 4.4.5.2 \textit{applied to $W_0$ instead of $W$}) \textit{if $W_0$ satisfies} Loc (7). \textit{Then the conditions} $\mathrm{A}(\cF_{W_0})$, $\mathrm{B}(\cF_{W_0})$, $\mathrm{C}(\cF_{W_0})$ (\textit{which we also denote simply by} $\mathrm{A}_{W_0}$, $\mathrm{B}_{W_0}$, $\mathrm{C}_{W_0}$) \textit{are equivalent.} (\textbf{\textit{N.B.}} \textit{If we do not make the preceding hypothesis on $W_0$, we will only have} $\mathrm{A}_{W_0} \Longrightarrow \mathrm{B}_{W_0} \Longrightarrow \mathrm{C}_{W_0}$, \textit{by} 4.4.2.)

4.4.7. We have now gone through all the conditions of ``fibration'' type on $W$, or on the pair $(W, W_0)$, which have seemed useful to me. I am going to lay them out here in a (more or less) summary diagram: % [sic] parues -> paru, agreement of the past participle

\msp{page 48}

\printedpage{28}
% head: A. Grothendieck : Dérivateurs, Ch. XVI, Édition des Universités de Montpellier II et Paris VII
% foot: Version 18 mars 2019

\begin{center}
\begin{tikzcd}[column sep=2.3em, row sep=small]
 & \overbrace{\mathrm{A}_{W_0}}^{\cF_{W_0},\,\Fibf_{W_0}} \arrow[dd, Rightarrow] & & \\
\overbrace{\mathrm{A}}^{\cF_W,\,\Fibf_W} \arrow[ur, Rightarrow] \arrow[dd, Rightarrow] & & & \\
 & \mathrm{B}_{W_0} \arrow[dd, Rightarrow] & & \\
\mathrm{B} \arrow[ur, Rightarrow] \arrow[dd, Rightarrow] & & & \\
 & \mathrm{C}_{W_0} \arrow[dr, Rightarrow, "\bullet"] \arrow[uuuu, Rightarrow, dotted, bend right=40, "\text{\scriptsize if } W_0\ \text{\scriptsize (7)}"'] & & \\
\mathrm{C} \arrow[ur, Rightarrow] \arrow[dr, Rightarrow, "\bullet"] & & \overbrace{\mathrm{D}_{W_0}}^{\Parf_{W_0}} \arrow[r, Rightarrow, "\bullet"] & \mathrm{F} : (\Cat, W, \Fib_W) \text{ a category with fibrations.} \\
 & \overbrace{\mathrm{D}}^{\Parf_W} \arrow[ur, Rightarrow] & &
\end{tikzcd}
\end{center}

\begin{center}
\begin{tikzcd}[row sep=large]
 & \mathrm{A} \arrow[r, Rightarrow] \arrow[d, Rightarrow] & \mathrm{A}_{W_0} \arrow[d, Rightarrow] & \\
\cF_W, \Fibf_W \smash{\left\{\vrule height 2.6cm depth 2.6cm width 0pt\right.} & \mathrm{B} \arrow[r, Rightarrow] \arrow[d, Rightarrow] & \mathrm{B}_{W_0} \arrow[d, Rightarrow] & \smash{\left.\vrule height 2.6cm depth 2.6cm width 0pt\right\}} \quad \text{\scriptsize\begin{tabular}{@{}l@{}}revised version\\ of the diagram\end{tabular}} \\
 & \mathrm{C} \arrow[r, Rightarrow] \arrow[d, Rightarrow] & \mathrm{C}_{W_0} \arrow[d, Rightarrow] & \\
 & \underbrace{\mathrm{D}}_{\Parf_W} \arrow[r, Rightarrow] & \underbrace{\mathrm{D}_{W_0}}_{\Parf_{W_0}} \arrow[r, Rightarrow] & \mathrm{F},
\end{tikzcd}
\end{center}

diagram which I am going to comment on. All the solid-line implications in the diagram have already been obtained or are tautological, except for the three $\overset{\bullet}{\Longrightarrow}$ to which I will return. The dotted implication $\mathrm{C}_{W_0} \dottedimplies \mathrm{A}_{W_0}$ is a conditional implication, subject to the condition marked on the diagram as $W_0$ (7), which means: $W_0$ satisfies Loc (7). This implication, \textit{i.e.} the equivalence of $\mathrm{A}_{W_0}$, $\mathrm{B}_{W_0}$, $\mathrm{C}_{W_0}$ under the condition $W_0$ (7), was seen in 4.4.6. I have marked, above the two groups (A, B, C) and ($\mathrm{A}_{W_0}$, $\mathrm{B}_{W_0}$, $\mathrm{C}_{W_0}$) as well as above D and D$'$ (which each correspond to such a group of conditions $\mathrm{A}(\cF)$, $\mathrm{B}(\cF)$, $\mathrm{C}(\cF)$ all equivalent, with $\cF \;=\; \Parf_W$ or $\cF = \Parf_{W_0}$) the
\msp{page 49}
sets of arrows which give rise to them respectively: it will be noted that there are two for the for the two triples named first, A, B, C being definable via $\cF_W$ or $\Fib_W$ indifferently, and likewise $\mathrm{A}_{W_0}$, $\mathrm{B}_{W_0}$, $\mathrm{C}_{W_0}$ via $\cF_{W_0}$ or $\Fib_{W_0}$, as already stated (4.4.3) and as we will see below (cor. 4.4.10, pages 54-55). On the other hand, the reduced triples D and $\mathrm{D}_{W_0}$ are each defined by just one $\cF$, namely by $\Parf_W$ and % [sic] pour les pour les -> pour les, words repeated; [sic] Fib_W / Fib_{W_0} -> Fibf_W / Fibf_{W_0}, the diagram and (4.4.3) give Fibf

\end{document}
